\section{Scheduling guarantees and limits of observation}\label{sec:scheduling}

\subsection{Comparisons that keep the baseline unchanged}\label{sec:effort-baseline}

A guarantee relative to the unchanged solver is available if the unchanged
solver runs unchanged. The following statements use an ideal preemptible work
model. Two solvers $A$ (the baseline) and $H$ (with the policy) run on the same
input with separate states and fixed random tapes. Each solver's sequence of
states depends only on the amount of work it has received, not on wall-clock
time or on the other solver. Run alone, they finish after $T_A$ and $T_H$ units
of work. A scheduler shares one processor; by aggregate time $t$, solver $S$
has received $a_S(t)$ units of work.

\begin{theorem}[Interleaving]\label{thm:interleaving}
Let $0<\alpha<1$ and $h\ge0$. Suppose that, as long as neither solver has
finished, $a_A(t)\ge(1-\alpha)t-h$ and $a_H(t)\ge\alpha t-h$. Then the first
of the two solvers finishes no later than
\begin{equation}\label{eq:interleaving}
 \min\left\{\frac{T_A+h}{1-\alpha},\ \frac{T_H+h}{\alpha}\right\}.
\end{equation}
\end{theorem}

\begin{proof}
Suppose neither solver has finished by $t_A=(T_A+h)/(1-\alpha)$. Then $A$ has
received at least $T_A$ units of work. Because its sequence of states depends
only on the work received, it has reached the state in which it finishes, a
contradiction. The same argument applies at $t_H=(T_H+h)/\alpha$.
\end{proof}

\begin{corollary}[Equal shares]\label{cor:equal-shares}
If the scheduler alternates quanta of length $q$ between the two solvers and
switching costs nothing, one of the solvers finishes by
$2\min\{T_A,T_H\}+q$.
\end{corollary}

\begin{proof}
Let $A$ receive the first quantum. During the period $[2kq,2(k+1)q]$, solver
$A$ holds at least $kq+\min\{t-2kq,q\}\ge t/2$ units of work and $H$ at least
$kq+\max\{t-2kq-q,0\}\ge t/2-q/2$. Apply Theorem~\ref{thm:interleaving} with
$\alpha=1/2$ and $h=q/2$: the bound is
$\min\{2T_A+q,\,2T_H+q\}$. Since $a_A(t)\ge t/2$, the term for $A$ can be
sharpened to $2T_A$.
\end{proof}

\begin{proposition}[Serial rescue]\label{prop:rescue}
Run $H$ for at most $b$ units of charged work, including its startup. If it
has not finished, discard it and run $A$ from its initial state with its
original input and random tape. The total work is $T_H$ if $T_H\le b$, and
$b+T_A$ otherwise. If the cap on $H$ can be overrun by at most $\delta$, the
second case becomes at most $b+\delta+T_A$.
\end{proposition}

\begin{proof}
If $H$ finishes within its allowance, the total is its own work. Otherwise its
charged work is at most $b$, or $b+\delta$ with overrun, and $A$ then executes
exactly its standalone trace, which costs $T_A$ including startup.
\end{proof}

Both guarantees depend on the baseline trace being unaffected. Shared memory or
cache contention, uncharged process startup, decisions that depend on
wall-clock time, and communication between the solvers can all violate the
premises. In particular, passing an incumbent or a row from $H$ to $A$ changes
$A$'s trajectory, and the bound $T_A$ no longer applies. The statements bound
work, not elapsed time on a shared machine. The computational comparison of
Section~\ref{sec:experiments} does not deploy either construction; it compares
independent full runs.

\subsection{History does not identify the best first action}\label{sec:effort-history}

The triggers and the pilot rule of Section~\ref{sec:algorithms-policy} choose
actions from what the policy has observed. The following example shows that
such a choice cannot be guaranteed correct from the observations alone.

\begin{example}[Indistinguishable histories]\label{ex:history}
Two actions each cost one unit. In environment~1 only action~1 produces a
proof; in environment~2 only action~2 does. Suppose the information available
to a selector before its first action is identical in the two environments.
A deterministic selector then makes the same first choice in both and chooses
wrongly in one of them. If a randomized selector chooses action~1 with
probability $p$, its first choice is correct with probability $p$ in
environment~1 and $1-p$ in environment~2, and $\min\{p,1-p\}\le1/2$.
\end{example}

The example limits only the first choice under identical information. It does
not rule out learning under statistical assumptions, structural tests that
distinguish the environments, or good average behavior on a class of
instances. Adaptive selection of solver components has been studied under such
assumptions, for example with bandit algorithms and online learning in
mixed-integer programming~\cite{hendel2018adaptive,chmiela2023scheduling}.
The point for this
paper is narrower. The triggers, ranking, and pilot of the measured policy are
heuristics; their value can be judged only by complete runs on instances, not
by their observed rewards.
